\chapter*{Introduction: One Screening, Several Films}
\addcontentsline{toc}{chapter}{Introduction: One Screening, Several Films}
\markboth{Introduction}{Introduction}

Two people sit next to each other in a cinema. The lights go down, the film begins, and both of them look at the same screen for the next hour and a half. When they leave, one of them has seen a comedy and the other a drama. Their films began together, reached their principal turns together, and ended together. The same characters moved through the same rooms. But the lines were delivered differently, the scenes were cut to a different rhythm, and the last ten minutes resolved in opposite directions. Neither person looked at a phone, wore a headset, or watched a second screen. Each wore a pair of glasses.

The scenario seems to require something impossible, because it seems to require that one illuminated surface present two different images at the same moment. It does not. A projected film is not a continuous image. It is a rapid sequence of discrete frames, each shown for a fraction of a second, and the impression of continuous motion is produced by the visual system rather than by the screen. Once the image is understood as a sequence, the scenario becomes a matter of scheduling. If the projector alternates frames from two films, and each viewer's glasses admit only the frames of one of them, then two films occupy one screen. Each viewer sees a complete film at half the frame rate the projector runs at, and neither sees the other's.

\section*{What is and is not new}

The optical principle is old and has been implemented several times. Laurens Hammond's Teleview system, patented in November 1922 (US 1,435,520) \cite{hammond1922} and first shown publicly at New York's Selwyn Theatre on December 27 of that year, gave each seat a rotary shutter synchronized to interlocked projectors alternating left-eye and right-eye images, so that each eye saw only its own frames and the audience saw depth. Teleview is a direct ancestor of the apparatus, though not of the idea: its two streams were halves of one picture, not two stories. Modern stereoscopic cinema and home television have used electronic shutter glasses for the same purpose. Sony's SimulView, announced in August 2011 for the PlayStation 3D Display released later that year \cite{simulview}, turned the stereoscopic arrangement sideways: instead of sending different frames to a viewer's two eyes, it sent different full-screen images to two players, so that each saw only their own view of a split-screen game without the split. It is the closest commercial precedent for co-located viewers receiving different full-screen streams, though its purpose was multiplayer play rather than authored films. Research on what has been called temporal psychovisual modulation has explored displays that show one image to viewers wearing synchronized glasses and a different one to the unaided eye \cite{wuzhai}.

None of this is claimed here as a discovery. What the book adds is what those systems did not need. A game shown to two players does not need the two views to tell a coherent story together. A stereoscopic film does not need its left and right images to diverge in meaning. The present proposal requires both. It asks for several films designed as members of one family: synchronized in time, sharing characters and events, differing in treatment, and able to be seen in the same room without any of them being the real one and the others alternatives. The engineering is a starting point. The subject is the form.

\section*{Neighbors and differences}

Several existing practices resemble the proposal in part, and the differences clarify what it is.

Stereoscopic cinema already sends different images to different eyes, but it recombines them into a single percept. The two streams are halves of one experience, not two experiences.

Captions and audio description add information to a common image without replacing it. A captioned viewer and an uncaptioned viewer see the same film with different supplements.

Alternate cuts, such as director's cuts, broadcast edits, and censored versions, are common, but they are distributed separately. They are watched at different times, in different places, by audiences who never share a room.

Branching film, in which the audience chooses what happens next, produces different experiences from one work, but only by interrupting it. \emph{Kinoautomat}, shown at Expo~67 in Montreal, stopped at decision points and let the audience vote, so that the whole room followed the majority's branch \cite{kinoautomat}: a shared choice with one outcome for everyone. Streaming branching video such as \emph{Black Mirror: Bandersnatch} gives each viewer or household its own path, but only by giving up the shared room \cite{bandersnatch}. Selective cinema lets each viewer take their own branch without stopping the film and without leaving the room.

Head-mounted displays can show every viewer something different, but they do so by replacing the shared screen with private ones. The room becomes a set of individuals who happen to be near each other.

Selective cinema, as the proposal will be called, keeps what each of these gives up. It keeps the common screen, the uninterrupted screening, and the shared room, and it still lets the audience see different films. The combination is what makes it more than a curiosity. It separates two things ordinary cinema has always bundled together: being present at the same event and receiving the same images.

\section*{The terms}

A few terms need fixing at the start.

A \emph{stream} is one sequence of frames within the projection, the frames a particular pair of glasses admits. A \emph{realization} is the complete audiovisual experience a viewer receives by following one stream, with its sound. A \emph{variant family} is the set of realizations designed together for synchronized presentation.

\emph{Blink-rate glasses} are glasses whose lenses open and close on a schedule synchronized to the projection. The name refers to the schedule of the shutter, not to the wearer's own blinking. The glasses do not watch the viewer; they are instruments of timing. A schedule is specified by two numbers, a period and a phase, and \Cref{ch:phase} shows that this pair, not a rate alone, is what selects a stream.

\emph{Selection} is choosing a stream. \emph{Concealment} is preventing someone from seeing a stream they have not chosen. The two are easily confused, and a large part of the book depends on keeping them apart. A mechanism that works well for selection may work poorly for concealment, and a system sold for concealment inherits obligations a system built for selection does not.

\section*{The argument}

Part~I treats the screen as a carrier. It asks how much temporal capacity a projector has, how many streams that capacity can support, what the glasses must do to separate them, and what the room sees when someone takes the glasses off. The answers impose costs that grow with every added stream, and they also reveal a design freedom that is easy to miss: the composite image seen without glasses can itself be designed.

Part~II treats the film as a family of realizations. Its central notion is \emph{switch admissibility}, the condition under which a viewer can move from one stream to another without the new stream presupposing something the viewer was never shown. Synchronization anchors, narrative rendezvous, and the accounting of time borrowed and repaid across streams all follow from that condition. The part then works through the principal kinds of variation: genre, explanatory depth, and endings.

Part~III is devoted to sound, which is harder to multiplex than light. Air cannot be interleaved the way frames can. The part defends a solution present in the idea from the start: a single composed score that serves every stream at once, keeping the room acoustically shared while the pictures diverge.

Parts~IV and~V take up viewpoint, local realization, and the practical work of writing, producing, and editing a variant family. Part~VI turns to the audience: what people share when they share a room but not a picture, what they say to each other afterward, and how and when they should be allowed to choose. Part~VII considers accessibility, instruction, and translation, where selective streams could do real good and could also sort people without their consent.

Part~VIII confronts the hardest questions. Someone must decide which versions exist and who can see them. The same glasses that let a viewer choose a comedy could let a distributor show different audiences different claims while everyone believes they saw one film. The same timing that selects a stream could be licensed so that only approved glasses can see at all. The part argues for disclosure, for glasses that record nothing, and for an account of the work's identity that locates it in its declared constraints rather than in any single sequence of frames. Part~IX sets out how such a system should be evaluated.

The principal claim of the book can be stated now:

\begin{quote}
Collective viewing does not require that everyone see the same thing. It requires that different experiences remain synchronized closely enough to count as one occasion.
\end{quote}

Everything that follows is an attempt to say what ``closely enough'' requires, what it costs, and who should decide.
